\chapter{Experimentální průzkum vlastností PFN}

V této kapitole provedeme řádu experimentu, ve kterých budeme zkoumat a ověřovat různé hypotézy o tom, jak funguje PFN.
Ačkoli Nagler se snaží ve své práci matematicky popsat principy PFN a naznačit jeho principy fungování, ale ve většině případu uvažuje jednovrstevný model a navíc vůbec nepopisuje, co PFN transformer ve výsledku vůbec dělá a jak to funguje uvnitř ho.
Proto se pokusíme tyto otevřené otázky probrat jestli ne skrz matematický aparát, tak alespoň pomocí empirických pozorování.
Všechny experimenty budou provdeny jenom pro architekturu typu transformer. 
Bude nás většinou zajímat vnitřní struktura PFN.
Rovněž prozkoumáme mechanismy, jak PFN vyhodnocuje vstupní data a čím se tyto mechanismy liší od regrese pomocí Gaussovských procesů.

\section{Vnitřní struktura PFN transformeru a jeho základní principy}

Nejdříve uvedeme základní analýzu PFN transformeru.
Zaměříme se zde na to, co přesně děla se vstupními daty, podíváme se, zde se snaží nakopírovat stejné mechanismy jako GP, anebo zda dělá něco jiného.
Experimenty jsme prováděli na dvou typech modelů.
Oba sdílí identickou architekturu transformeru: $6$ vrstev, $8$ attention hlav, dimenze embeddingu $512$, dimenze skrytých vrstev dopředné sítě $1024$ a výstupní \texttt{FullSupportBarDistribution} s $1000$ biny, celkem $14{,}14$ milionu parametrů.
Architektura je nastavena s parametry \lstinline|features_per_group=1| a \lstinline|attention_between_features=False|, tedy přesně tak, jak jsme ji popsali v předchozí kapitole pro případ $1$D regrese.
Oba modely byly trénovány na syntetických datech vygenerovaných z Gaussovského procesu s RBF kernelem, vstupní prostor $x \in [0, 1]$.
Liší se však tím, z jakého prioru nad hyperparametry trénovací data pocházela, a z toho plynoucí délkou tréninku.

\textbf{Model se fixními hyperparametry:} 

(\textit{100-epoch model}) byl trénován na GP s pevně nastavenými hyperparametry RBF kernelu: $\ell = 0{,}3$, $\text{outputscale} = 1{,}0$, $\sigma^2 = 10^{-4}$.
Model se tedy učil aproximovat posterior jednoho konkrétního GP a trénink konvergoval za $100$ epoch.

\textbf{Model s náhodně vzorkovanými hyperparametry:} 

(\textit{random-HP model}) byl trénován na GP, jehož hyperparametry byly při každé dávce vzorkovány nezávisle z předem dané distribuce: $\ell \sim \mathcal{U}(0{,}05,\, 1{,}0)$, $\text{outputscale} \sim \text{LogNormal}(0,\, 0{,}5)$, $\sigma^2 \sim 10^{\mathcal{U}(-5,\,-2)}$.
Tento model se tedy musí naučit generalizovat přes celou rodinu GP funkcí s různými korelačními délkami a amplitudami — jde o plnohodnotný meta-learning nad distribucí hyperparametrů.
Právě proto vyžaduje podstatně delší trénink: model viděl data z nespočetně různých tvarů funkcí a musel se naučit rozpoznávat hyperparametry přímo z kontextových dat.
Trénink trval $500$ epoch.
\par

Výsledky obou modelů porovnáváme v průběhu jednotlivých experimentů mezi sebou a s referenčním analytickým GP posteriorem, který pro dané hyperparametry a kontextová data vrací přesné řešení.
Začneme ovšem s případem fixních hyperparametrů. 
Pro model to je značné zjednodušení, přesně ví, co se má naučit, a jak se ukáže dále, velmi efektivně.
Na tomto jednoduchém příkladě ukážeme hlavní principy inference této metody, abychom následně mohli přejít na těžší příklady.
\par

\subsection{Struktura attention matic}

Začneme však s tím, že vykreslíme si jednotlivé vrstvy a podíváme se na vnitřnosti transformeru a jak se jeho vrstvami prochází tok dat.
Attention matice je přirozené okno do toho, jak transformer interně organizuje informaci.
Připomínáme, že v naší konfiguraci má model $6$ vrstev a v každé vrstvě $8$ attention hlav, přičemž každá hlava může sledovat jiný aspekt vstupních dat.
Vizualizaci provedeme pro 100-epoch model na sekvenci délky $120$: $20$ trénovacích bodů a $100$ testovacích bodů.
Výsledkem je attention matice tvaru $120 \times 120$, přičemž každý prvek $a_{ij}$ říká, jak moc bod $i$ (jako \textit{query}) věnuje pozornost bodu $j$ (jako \textit{key}).
Matici průměrujeme přes $8$ attention hlav dané vrstvy, aby byl celkový obraz přehledný.
\par

Protože víme, které pozice odpovídají trénovacím bodům ($0$–$19$) a které testovacím ($20$–$119$), lze matici přirozeně rozdělit na čtyři kvadranty:

\begin{center}
\begin{tabular}{c|c|c}
 & \textbf{Key: Train} & \textbf{Key: Test} \\
\hline
\textbf{Query: Train} & Train$\to$Train & Train$\to$Test \\
\hline
\textbf{Query: Test}  & Test$\to$Train  & Test$\to$Test  \\
\end{tabular}
\end{center}

Kvadrant \textit{Test$\to$Train} je pro nás klíčový: říká, jak moc testovací body ``čerpají'' informaci z trénovacích dat při budování své predikce.
Z pohledu GP inference bychom očekávali, že právě tento kvadrant bude dominantní — predikce v novém bodě $x^*$ závisí na pozorovaných párech $(X, y)$, nikoli naopak.
\par

Na obrázku~\ref{fig:attention_all_layers} je vidět, jak se struktura attention matic vyvíjí přes vrstvy:

\textbf{Vrstva 0:}

vykazuje relativně rovnoměrné, distribuované attention váhy.
Žádný kvadrant výrazně nedominuje a váhy jsou rozloženy plošně přes mnoho pozic.
Tuto vrstvu interpretujeme jako \textit{explorativní} fázi: model sbírá globální informaci o rozložení všech bodů v sekvenci bez toho, aby se specializoval.

\textbf{Vrstvy 1–4:}

attention postupně řídne — váhy se soustřeďují na méně pozic a zbytek klesá k nule.
Vzory se stávají strukturovanějšími: kvadrant \textit{Test$\to$Train} začíná nabývat na intenzitě, zatímco \textit{Train$\to$Test} slábne.
Tyto prostřední vrstvy zřejmě slouží k internímu zpracování a redistribuci informace — postupnému sestavování reprezentace, ze které bude moci poslední vrstva provést samotnou inferenci.

\textbf{Vrstva 5:}
 
(poslední) vykazuje nejzřetelnější vzor.
Kvadrant \textit{Test$\to$Train} je jasně dominantní a strukturovaný, zatímco kvadrant \textit{Train$\to$Test} je téměř nulový.
Tato asymetrie je přesně tím, co bychom čekali od GP inference: testovací body jsou podmíněny trénovacími daty, nikoli naopak.
\par

Podstatné je, že tuto kauzální asymetrii model \textbf{neměl explicitně zakódovanou} — vyplynula přirozeně z optimalizace na syntetických GP datech.
Kdybychom totiž model trénovali na datech bez kauzální struktury, attention by zůstala symetrická.
Samotná existence asymetrie v poslední vrstvě je tak prvním experimentálním důkazem, že PFN provádí Bayesovskou inferenci — a nikoliv pouhý pattern matching nebo memorování trénovacích příkladů.

\begin{figure}[p]
    \centering
    \includegraphics[width=\linewidth]{figures/fig_big_03_attention_all_layers.png}
    \caption{Attention matice pro všech $6$ vrstev 100-epoch modelu (průměr přes $8$ hlav) pro jednu konkrétní GP realizaci. Osa $x$: \textit{key} pozice (na které body se díváme), osa $y$: \textit{query} pozice (které body se dívají). Cyánová čára odděluje trénovací ($0$–$19$) a testovací ($20$–$119$) pozice. Světlejší barva odpovídá vyšší attention váze. Přesné hodnoty vah závisí na konkrétních vstupních datech, kvalitativní vzor (vrstva 0 distribuovaná $\to$ vrstva 5 s dominantním \textit{Test$\to$Train} blokem) je však stabilní.}
    \label{fig:attention_all_layers}
\end{figure}

\clearpage
\subsection{Matematický rámec pro srovnávací experimenty}

Než přistoupíme k samotným experimentům, je užitečné shrnout matematický aparát, který budeme v dalších sekcích potřebovat.
Jde konkrétně o tři formule, jejichž vzájemný vztah je jádrem naší analýzy.

\textbf{Posteriorní střední hodnota GP.}

Po pozorování kontextových párů $(X, y)$ je posteriorní střední hodnota Gaussovského procesu v novém bodě $x^*$ dána vztahem:
$$\mu(x^*) = k(x^*, X)\bigl[K(X,X) + \sigma^2 I\bigr]^{-1} y,$$
kde $K(X,X)_{ij} = k(x_i, x_j)$ je kernelová matice a $\sigma^2$ je variance šumu.
Klíčovou roli hraje inverze $[K + \sigma^2 I]^{-1}$, která \textit{dekore\-luje} blízké trénovací body — body, které si jsou podobné, si vzájemně ``konkurují'' a jejich kumulativní vliv je potlačen.

\textbf{Nadaraya-Watsonův estimátor.}

Nejjednodušší alternativou je Nadaraya-Watsonův (NW) estimátor, tj. normalizovaný váhovaný průměr trénovacích hodnot:
$$\hat{y}_{\text{NW}}(x^*) = \frac{\displaystyle\sum_i k(x^*, x_i)\, y_i}{\displaystyle\sum_i k(x^*, x_i)}.$$
NW se od GP liší právě absencí inverze kernelové matice — implicitně předpokládá, že trénovací body jsou vzájemně nekorelované.
Blízké body proto dostávají neúměrně velkou kumulativní váhu a predikce je příliš hladká.

\textbf{Attention mechanismus transformeru.}

Attention v PFN počítá pro každý testovací bod $x^*$ váhovanou kombinaci hodnot z trénovacích bodů:
$$\text{Attention}(Q, K, V) = \text{softmax}\!\left(\frac{QK^\top}{\sqrt{d_k}}\right)V, \quad Q = XW_Q,\; K = XW_K,\; V = XW_V.$$
Výsledná kernelová matice je \textit{datově závislá}:
$$\text{Kernel}_{\text{attn}}(X) = \text{softmax}\!\left(\frac{X W_Q W_K^\top X^\top}{\sqrt{d_k}}\right).$$
Na rozdíl od fixního RBF kernelu se tato matice plně adaptuje na konkrétní vstupní sekvenci.
Matice vah $W_Q, W_K$ jsou naučeny z dat a mohou kódovat libovolnou strategii podobnosti.

Klíčová otázka, kterou budeme v následujících experimentech testovat, pak zní: \textit{odpovídají attention váhy PFN normalizovaným RBF vahám, nebo se model naučil kvalitativně odlišnou strategii?}
A pokud odlišnou — jak moc se jeho výstup blíží skutečnému GP posterioru oproti triviálnímu NW estimátoru?

\subsection{Dělá PFN něco podobného jako GP s RBF kernelem?}

V této sekci se zaměříme, zda se PFN transformer snaží napodobit chování RBF kernelu, a proto ho tak dobře aproximuje, anebo dělá něco jiného.
Konkrétně pro každý testovací bod $x^*$ porovnáme attention váhy $a_i$ z poslední vrstvy (průměr přes $8$ hlav) s normalizovanými RBF vahami:
$$w_i = \frac{k(x^*, x_i)}{\sum_j k(x^*, x_j)}.$$
Experiment provedeme pro 100-epoch model na $n_{\text{ctx}} = 20$ trénovacích bodech a $100$ testovacích bodech, průměrem přes $5$ nezávislých seedů.
\par

Každá metrika je průměrována přes $5$ nezávislých seedů, přičemž pro každý seed je vygenerována jedna GP sekvence délky $100$ (prvních $20$ bodů jako kontext, všech $100$ jako testovací body).
MSE je počítáno jako střední kvadratická odchylka mezi střední hodnotou predikce PFN a analytickým GP posteriorem přes všech $100$ testovacích bodů.
Korelace attention vs.\ RBF je pro každý testovací bod $x^*$ spočítána jako Pearsonova korelace mezi attention vahami poslední vrstvy (průměr přes $8$ hlav) a normalizovanými RBF vahami $\exp(-d^2/2\ell^2)$, zprůměrovaná přes všechny testovací body.
Naměřené hodnoty jsou shrnuty v tabulce~\ref{tab:exp3}.

\begin{table}[htbp]
\centering
\begin{tabular}{lc}
\hline
\textbf{Metrika} & \textbf{100-epoch model} \\
\hline
Průměrná korelace Attn vs RBF & $0{,}374 \pm 0{,}270$ \\
MSE(PFN, GP)  & $0{,}000167 \pm 0{,}000096$ \\
MSE(NW, GP)   & $0{,}2099 \pm 0{,}1462$ \\
MSE(PFN, NW)  & $0{,}2092 \pm 0{,}1439$ \\
Poměr MSE(NW)/MSE(PFN) & ${\sim}1\,256\times$ \\
\hline
\end{tabular}
\caption{Výsledky korelační analýzy attention vah vs.\ RBF kernelu a srovnání přesnosti predikcí. Průměr přes $5$ seedů.}
\label{tab:exp3}
\end{table}

\begin{figure}[htbp]
    \centering
    \subcaptionbox{Detail attention vah poslední vrstvy (průměr přes $8$ hlav) pro vybraný testovací bod v porovnání s normalizovanými RBF vahami. Attention je výrazně ostřejší — koncentruje téměř veškerou váhu na nejbližší bod — zatímco RBF distribuuje váhu plynule. Korelace $0{,}374 \pm 0{,}270$.\label{fig:attn_vs_rbf}}[1.0\linewidth]{%
        \includegraphics[width=0.82\linewidth]{figures/fig_big_04_attention_detail.png}}
    \vspace{0.5em}
    \subcaptionbox{Srovnání predikcí PFN, Nadaraya-Watsonova estimátoru a analytického GP posterioru ($n_{\text{ctx}} = 20$). NW systematicky přehlazuje, PFN sleduje GP posterior s MSE $167 \times 10^{-6}$ oproti MSE $0{,}21$ u NW — rozdíl ${\sim}1\,256\times$.\label{fig:pfn_vs_nw}}[1.0\linewidth]{%
        \includegraphics[width=0.82\linewidth]{figures/fig_big_05_nw_comparison.png}}
    \caption{100-epoch model: porovnání attention vah s RBF kernelem (nahoře) a predikcí PFN, NW a GP (dole).}
\end{figure}
\par

Z tabulky~\ref{tab:exp3} vyplývají dvě zásadní pozorování.
Za prvé, korelace attention vah s RBF kernelem je nízká ($0{,}374 \pm 0{,}270$) a s vysokou variabilitou, což potvrzuje, že model \textit{neimplementuje} RBF kernel-smoothing.
Attention se naučila kvalitativně odlišnou strategii: jak je vidět na obrázku~\ref{fig:attn_vs_rbf}, attention váhy jsou výrazně ostřejší než RBF — soustřeďují téměř veškerou váhu na nejbližší bod, zatímco RBF ji distribuuje plynule.

Za druhé, MSE(PFN, NW) $\approx$ MSE(NW, GP) — to znamená, že NW a PFN predikují \textit{zcela odlišné věci}: PFN je o $1\,256\times$ přesnější než NW.
Na obrázku~\ref{fig:pfn_vs_nw} je vidět, že NW systematicky přehlazuje oblasti s hustými trénovacími body, zatímco PFN posterior sleduje analytické GP řešení.
\par

Proč je attention ostřejší než RBF? Attention hlavy nepotřebují počítat čistou kernelovou podobnost $k(x^*, x_i)$.
Potřebují vypočítat to, co v kombinaci s ostatními hlavami a FFN vrstvami dá nejlepší aproximaci GP posterioru.
A pro ten může být výhodnější ostřejší strategie — FFN pak snáze provede korekci odpovídající efektu $K^{-1}$.
\par

\textbf{Epistemologická poznámka: co korelace Attn vs RBF skutečně říká.}

Nízká korelace attention vah s RBF kernelem je platný a hodnotný závěr o chování \textit{jedné komponenty} sítě.
Existují však důvody, proč z ní nelze přímo usuzovat na celkové chování modelu.

Za prvé, attention není výstupní mechanismus — po ní následuje FFN, residuální spojení a v případě vícevrstvého modelu dalších pět bloků.
I kdyby attention váhy v poslední vrstvě přesně odpovídaly normalizovaným RBF vahám, výstup PFN by se od NW stále lišil díky korekcím v FFN.

Za druhé, správnou referenční hodnotou pro porovnání \textit{nejsou} čisté RBF váhy $k(x^*, x_i)$, ale efektivní GP váhy:
$$w_i^{\text{GP}} = \sum_j k(x^*, x_j)\,(K^{-1})_{ji}.$$
Tyto váhy jsou nelokální: blízké trénovací body si vzájemně konkurují a po aplikaci $K^{-1}$ mohou mít menší váhu, než by odpovídalo RBF, zatímco vzdálené body mohou dostat kladnou či dokonce zápornou váhu.
Nízká korelace attention s čistým RBF je tedy strukturálně očekávána — i perfektní GP aproximátor by vykazoval nízkou korelaci s RBF při tomto experimentálním designu.

Shrnutí: nízká korelace Attn vs RBF potvrzuje, že PFN neprovádí prostý kernel averaging.
Kvalitu celkových predikcí vůči GP pak jednoznačně zachycuje poměr MSE ${\sim}1\,256\times$ — a ten hovoří zcela ve prospěch PFN.

\subsection{Jak přesný je PFN na skutečné GP realizaci?}

Předchozí experiment porovnával PFN se syntetickým GP posteriorem — modelu jsme předložili trénovací body a srovnávali jeho predikci s analytickým řešením.
Nyní přistoupíme k přísnějšímu testu: oba modely (PFN i GP oracle) budou predikovat \textit{skutečné šumové hodnoty} neviděné GP realizace.
Cílem je zjistit, jak se přesnost PFN vyvíjí s velikostí kontextu $n_{\text{ctx}} \in \{5, 10, 15, 20\}$.
\par

Vygenerujeme kompletní GP realizaci délky $100$ bodů.
Z ní náhodně vybereme $n_{\text{ctx}}$ bodů jako trénovací kontext; zbývající body slouží jako testovací cíl.
MSE počítáme jako střední kvadratickou odchylku predikce od skutečných (šumových) $y$-hodnot přes testovací body.
GP oracle zná správné hyperparametry $(\ell = 0{,}3,\; \sigma^2 = 10^{-4})$ a počítá analytický posterior.
Metriky průměrujeme přes $5$ seedů $\times$ $20$ realizací $=$ $100$ instancí na každé $n_{\text{ctx}}$.
\par

\begin{table}[htbp]
\centering
\begin{tabular}{lcc}
\hline
$n_{\text{ctx}}$ & \textbf{PFN (100-epoch)} & \textbf{GP oracle} \\
\hline
$5$  & $0{,}0408 \pm 0{,}0698$ & $\mathbf{0{,}0340 \pm 0{,}0635}$ \\
$10$ & $0{,}0123 \pm 0{,}0349$ & $\mathbf{0{,}0112 \pm 0{,}0360}$ \\
$15$ & $0{,}0064 \pm 0{,}0271$ & $\mathbf{0{,}0028 \pm 0{,}0036}$ \\
$20$ & $0{,}0021 \pm 0{,}0025$ & $\mathbf{0{,}0017 \pm 0{,}0011}$ \\
\hline
\end{tabular}
\caption{MSE vůči skutečným hodnotám GP realizace jako funkce velikosti kontextu. Průměr přes $100$ instancí ($5$ seedů $\times$ $20$ realizací).}
\label{tab:exp4}
\end{table}

\begin{figure}[htbp]
    \centering
    \subcaptionbox{MSE PFN (modrá) a GP oracle (zelená přerušovaná) jako funkce $n_{\text{ctx}}$. Chybové úsečky $= \pm 1$ std přes $20$ realizací. GP oracle mírně vede od $n = 10$, protože zná správné $\ell$.\label{fig:exp4_mse}}[0.82\linewidth]{%
        \includegraphics[width=0.82\linewidth]{figures/fig_big_06_context_size_mse.png}}
    \vspace{0.5em}
    \subcaptionbox{Vizualizace predikcí pro $n \in \{5, 10, 15, 20\}$. Červené body: kontextová data dostupná modelu; šedé body: skutečné neviděné hodnoty; modrá: PFN; zelená přerušovaná: GP posterior se správným $\ell$. Všechny řádky vychází ze stejné GP realizace, liší se pouze velikostí kontextu. Od $n = 15$ obě křivky téměř splývají.\label{fig:exp4_grid}}[0.82\linewidth]{%
        \includegraphics[width=0.82\linewidth]{figures/fig_big_06b_context_size_grid.png}}
    \caption{100-epoch model: MSE a predikce jako funkce velikosti kontextu.}
    \label{fig:exp4_combined}
\end{figure}

\textbf{Diskuze:}

Z tabulky~\ref{tab:exp4} a obrázku~\ref{fig:exp4_combined} vyplývají tři pozorování.

Za prvé, GP oracle konzistentně dosahuje nižšího nebo srovnatelného MSE oproti PFN ve všech hodnotách $n_{\text{ctx}}$.
To je očekávané: oracle zná přesné hyperparametry a provádí exaktní Bayesovskou inferenci, zatímco PFN je pouze aproximace.

Za druhé, rozdíl se zvětšuje se středním $n_{\text{ctx}}$.
Při $n = 5$ je převaha GP oracle malá ($0{,}034$ vs $0{,}041$); při $n = 15$ již GP dosahuje $0{,}003$ oproti $0{,}006$ u PFN.
Příčinou je, že s více kontextovými body disponuje GP oracle přesnou inversí kernelové matice $K^{-1}$, zatímco PFN tuto operaci pouze aproximuje.
Při malém kontextu je $K^{-1}$ málo informativní a obě metody jsou si blíže.

Za třetí, od $n = 20$ se obě metody přibližují: MSE($0{,}0021$ vs $0{,}0017$) jsou si prakticky rovnocenné.
PFN tedy s dostatečným kontextem konverguje k chování GP oracle, aniž by explicitně invertoval kernelovou matici.
Toto je jedním z klíčových empirických potvrzení toho, že PFN provádí plnou GP inferenci — nikoliv pouze kernel smoothing.

\subsection{Konvergence střední hodnoty a kalibrace variance}

Experiment 1 zkoumá, zda PFN správně předpovídá nejen střední hodnotu, ale také tvar profilu uncertainty — tedy kde je model více a kde méně jistý.
Trénovací data jsou omezena na region $x \in [0{,}3,\, 0{,}7]$ ($n_{\text{ctx}} = 20$), testovací body pokrývají celý interval $[0, 1]$ ($n_{\text{test}} = 100$).
Průměr přes $5$ realizací $\times$ $5$ seedů $= 25$ instancí celkem.
Metriky kalibrace: $\text{Corr}(\hat{\sigma}_{\text{PFN}},\, \sigma_{\text{GP}})$ (korelace směrodatných odchylek) a $\text{MSE}(\hat{\sigma}_{\text{PFN}},\, \sigma_{\text{GP}})$ (absolutní chyba kalibrace).
Dále uvádíme \textit{ratio far/near} — poměr průměrné std v regionu extrapolace ($x \in [0,\, 0{,}2] \cup [0{,}8,\, 1]$) k std v centru dat ($x \in [0{,}4,\, 0{,}6]$), který měří, jak výrazně model zvyšuje nejistotu mimo region trénovacích dat.
\par

Střední hodnota PFN sleduje GP posterior s vysokou přesností v regionu s daty.
Mimo trénovací region ($x < 0{,}3$, $x > 0{,}8$) PFN správně konverguje k prioru (hodnota~$0$), stejně jako GP.
Ojediněle, pro realizace s vysokou amplitudou ($|m_{\text{train}}| > 1{,}5$), dochází k neúplné konvergenci — GP se v takových případech chová identicky.

\begin{table}[htbp]
\centering
\begin{tabular}{lc}
\hline
\textbf{Metrika} & \textbf{Hodnota} \\
\hline
Corr(PFN std, GP std) & $\mathbf{0{,}9812 \pm 0{,}0252}$ \\
MSE(PFN std, GP std) & $0{,}00376 \pm 0{,}00911$ \\
Ratio far/near — PFN & $\mathbf{21{,}3\times \pm 5{,}2\times}$ \\
Ratio far/near — GP  & $25{,}9\times \pm 3{,}8\times$ \\
\hline
\end{tabular}
\caption{Experiment 1: kalibrace variance pro 100-epoch model (průměr přes 25 instancí).}
\label{tab:exp1}
\end{table}

\begin{figure}[htbp]
    \centering
    \subcaptionbox{Srovnání střední hodnoty a variance PFN (modrá) a GP (zelená). Trénovací body v $x \in [0{,}3,\, 0{,}7]$, predikce v celém $[0, 1]$. PFN správně rozpoznává region extrapolace zvýšenou variancí.\label{fig:exp1_uncertainty}}[0.82\linewidth]{%
        \includegraphics[width=0.82\linewidth]{figures/fig_big_02_uncertainty.png}}
    \caption{Experiment 1 — 100-epoch model: konvergence střední hodnoty a kalibrace variance.}
    \label{fig:exp1_combined}
\end{figure}

\textbf{Diskuze:}

Korelace $r = 0{,}98$ potvrzuje, že PFN správně identifikuje \textit{tvar} profilu uncertainty — kde je nejistota větší a kde menší.
$\text{MSE}(\hat{\sigma}) = 0{,}0038$ signalizuje, že absolutní kalibrace je přesná.
Ratio far/near u PFN ($21{,}3\times$) dosahuje přibližně $82\,\%$ dynamického rozsahu GP ($25{,}9\times$): PFN tedy mírně podhodnocuje míru zvýšení uncertainty v regionu extrapolace, avšak kvalitativně chování zachovává.
Toto je zásadní vlastnost: model bez explicitní representace GP posterioru přesto generuje uncertainty, která odpovídá bayesovskému výpočtu.
